\chapter{Conclusion}
\label{chap:conclusion}

\section{Summary of the project}

This semester project started from an existing implementation of the Sponsored Fair
Exchange protocol and turned it into a more structured, measurable, and gas-aware research
prototype. The goal was not to redesign the cryptographic protocol from first principles.
The goal was to understand the inherited implementation, identify which protocol variants
matter for gas consumption, implement the most relevant ones, and measure them under
clearly defined scopes.

The first part of the work was therefore investigative. The existing system combined Rust
and WebAssembly precontract generation, Solidity optimistic and dispute contracts, a
Next.js/Electron interface, a SQLite backend, and an ERC-4337-style bundler path. Before
changing the protocol logic, it was necessary to reconstruct how these components actually
interacted. This included rebuilding the local stack, mapping UI and contract actions to
the steps of the protocol, and identifying where deployment, sponsorship, dispute
triggering, and proof generation really occurred in the code.

The second part of the work focused on variants. The project implemented and tested the
main optimistic simplifications: no \(S\)-deposit, \(S=B\), and \(S=V\). It also implemented
the relevant self-sponsored dispute behavior for \(S_B=B\) and \(S_V=V\), and a dedicated
hardcoded \(\mathsf{SHA256}\) dispute mode for \(desc=\mathsf{SHA256}(x)\). These variants
were first made functional in a monolithic implementation, then reorganized into a final
specialized architecture with smaller contracts and clone-based optimistic deployments.

The third part of the work was experimental. The project produced same-scope gas
measurements for the initial implementation, the monolithic intermediate implementation,
and the final specialized implementation. It also measured off-chain runtime for large-file
precontract and dispute witness generation. The most important large-file result is that a
complete hardcoded \(\mathsf{SHA256}\) dispute was executed for a 1 GiB input using the
native Rust pipeline.

\section{Summary of results}

The main positive result is on the optimistic path. In the initial implementation, an
honest optimistic exchange cost \(2,300,201\) gas, including deployment and execution. In
the final clone-based specialized architecture, the marginal optimistic cost per exchange
falls to \(500,720\) gas for \texttt{no\_S\_deposit}, \(516,349\) gas for \(S=B\), and
\(556,309\) gas for \(S=V\), once the reusable infrastructure has been deployed. This is the
clearest gas reduction obtained in the project.

The second result is methodological but important: the intermediate monolithic implementation
was not an optimization. It was a functional intermediate step. It made the variants
executable, but it increased gas because the deployed bytecode contained many mutually
exclusive branches. This negative result directly motivated the final split architecture.
It also made clear that gas comparisons must distinguish one-time infrastructure costs,
marginal per-exchange costs, dispute deployment, post-trigger dispute execution, and local
Step 8 costs.

The third result is the dispute-side specialization. On the same 4 MiB dispute
initialization scope, the final specialized implementation costs \(5,495,184\) gas, compared
with \(7,054,365\) gas for the initial implementation and \(8,872,096\) gas for the
monolithic intermediate implementation. On the 16 KiB same-scope dispute comparison, all
four final specialized variants improve over their monolithic counterparts, with the
self-sponsored hardcoded case reaching \(4,042,185\) gas in the measured scope.

The fourth result is the hardcoded \(\mathsf{SHA256}\) mode. In this mode, the selected
circuit is fixed by the precontract metadata. The dispute contract reconstructs the
relevant gate in Step 8 and no longer requires the generic circuit membership proof
\(\pi_1\). Locally, this saves \(62,542\) gas on the measured 16 KiB Step 8 call and about
\(166,000\)--\(169,000\) gas on the isolated 1 GiB \(h_{\mathit{circuit}}\) check. In the
full specialized implementation, the hardcoded path also reduces the dispute deployment
component because it uses a dedicated deployer and dispute account.

The fifth result is practical scalability. The native Rust pipeline can handle the 1 GiB
hardcoded benchmark. The measured self-sponsored hardcoded path executes the optimistic
setup, the dispute trigger, 26 challenge rounds, Step 8, and finalization for a total of
\(13,990,092\) gas. This is a concrete end-to-end result for the implemented benchmark
path. It should be interpreted correctly: it is a native CLI benchmark, not a
browser-only workflow.

\section{Final assessment}

The project achieved its main implementation and measurement objectives. The inherited
system was made executable and analyzed. The main sponsoring variants were specified,
implemented, and tested. The monolithic approach was identified as gas-inefficient and
replaced by a specialized architecture. The optimistic path was substantially reduced in
marginal gas cost. The hardcoded \(\mathsf{SHA256}\) dispute mode was implemented with the
contract reconstructing the gate instead of relying on a vendor-provided generic gate and
\(\pi_1\). Finally, the project demonstrated a 1 GiB dispute path using native tooling.

At the same time, the result remains a research prototype. The implementation still uses a
local development UI, local accounts, a local backend, and Hardhat measurements. The
optimized clone paths do not provide a full production wallet and ERC-4337 experience. The
precontract verification route exists, but acceptance is not yet enforced by a backend
verification precondition. The proof-generation API is also partial. These limitations are
now documented explicitly rather than hidden behind the gas results.

The project also identified inherited cryptographic-format issues. The generic V2 circuit
encoding is functional for the generated circuits used in the tests, but it is not the
cleanest final specification because some gate arities are inferred rather than encoded
explicitly. The Merkle-style accumulators also lack domain separation between leaves and
internal nodes. These issues were not introduced by the new variants, but they should be
fixed in a future protocol version before any production deployment or paper-ready
specification.

Overall, the final implementation is stronger than the starting point in three ways. It is
more gas-efficient on the main same-scope measurements, it is more explicit about which
protocol variants are being executed, and it is more honest about its remaining security
and engineering boundaries. This combination is important: a useful research prototype
should not only produce lower numbers, but also make clear what those numbers mean.

\section{Perspectives}

The natural next step is not to add more flags to the existing contracts. The project
showed that this path leads back to monolithic bytecode and confusing measurements. Future
work should instead continue the specialization and deployment-architecture direction.

A first direction is a cleaner factory-first or registry-first design. Such an architecture
could represent exchanges by identifiers or deterministic minimal proxies, reuse deployed
logic more systematically, and make Step 1+2 fusion possible even when \(S \neq B\). This
would require a careful redesign of signatures, state layout, and ERC-4337 validation, but
it is the most coherent path for reducing repeated deployment costs further.

A second direction is a new commitment format. A V3 circuit encoding should remove
ambiguity by encoding arity explicitly or by using distinct opcodes for cases such as
one-input versus two-input SHA-256 gates and zero-input versus one-input CONST gates. The
Merkle accumulators should also be domain-separated, for example by using different hash
prefixes for leaves and internal nodes. This would make the implementation cleaner from a
cryptographic engineering perspective.

A third direction is completing the verification workflow. In a production-quality system,
the buyer should not be able to accept a precontract unless the relevant verification has
succeeded and the verified values are bound to the exchange. Similarly, the dispute UI
should make actor-side verification explicit before each signed action. This would turn the
existing verification tools into enforced protocol workflow steps.

A fourth direction is optimizing the remaining dispute bottlenecks. After the optimistic
path improvements, the dominant costs are dispute deployment, \texttt{triggerDispute}, and
Step 8 evaluation, especially for AES-related gates and the associated
\(h_{\mathit{ct}}\)/\(hpre\) proofs. Further work should target these components directly.

Finally, the 1 GiB benchmark suggests a practical architecture for large-file fair exchange:
the frontend should orchestrate the user workflow, while native backend workers generate
large cryptographic artifacts. Turning this into a robust user-facing system would require
streaming, progress reporting, artifact availability, wallet integration, and production
deployment assumptions. The present project provides the measurements and implementation
foundation needed to make that next step precise.
